% !PRINTED 159
\weberpage{159}

In the region of binary substitutions, the substitution $A$:
\[
(y_1, y_2) = \begin{pmatrix} a, & b \\ c, & d \end{pmatrix} (x_1, x_2)
\]
regarded as a substitution of the ratios, is equivalent to the \gesperrt{linear fractional substitution}:
\[
\eta = \frac{a\,\xi + b}{c\,\xi + d},
\]
if we set $\xi = x_1 : x_2$ and $\eta = y_1 : y_2$; if a second substitution $A'$:
\[
(z_1, z_2) = \begin{pmatrix} a', & b' \\ c', & d' \end{pmatrix} (y_1, y_2),
\]
or
\[
\zeta = \frac{a'\eta + b'}{c'\eta + d'}
\]
is added, one obtains the composite substitution
\[
A'' = A'A,
\]
by which $\zeta$ is expressed in terms of $\xi$:
\[
\zeta = \frac{a''\xi + b''}{c''\xi + d''},
\]
according to the rules of composition, in the form
\[
\begin{pmatrix} a'', & b'' \\ c'', & d'' \end{pmatrix} = \begin{pmatrix} a', & b' \\ c', & d' \end{pmatrix} \begin{pmatrix} a, & b \\ c, & d \end{pmatrix} = \begin{pmatrix} a'a + b'c, & a'b + b'd \\ c'a + d'c, & c'b + d'd \end{pmatrix}.
\]
If such a substitution is multiplicative, that is, if it has the form
\[
\begin{pmatrix} a, & 0 \\ 0, & d \end{pmatrix},
\]
then we shall also call the ratio $a : d$ the \gesperrt{multiplier}.

\begin{center}
\S.\,39.

\gesperrt{Permutations as linear substitutions.}
\end{center}

The importance of the linear substitutions, and especially of the \gesperrt{finite} groups formed from them, for algebra results from the fact that the permutation groups of $n$ elements can be regarded as special cases of such substitution groups.
